\section{Постановка задачи}
\label{sec:problem}
\begingroup
\raggedright

Структура данных по варианту: PATRICIA. Необходимо создать программную библиотеку, реализующую указанную структуру данных, на основе которой разработать программу-словарь. В словаре каждому ключу, представляющему собой регистронезависимую последовательность букв английского алфавита длиной не более 256 символов, поставлен в соответствие некоторый номер от \(0\) до \(2^{64}-1\). Разным словам может быть поставлен в соответствие один и тот же номер.

Ввод: стандартный ввод или \texttt{input.txt}. Вывод: стандартный вывод или \texttt{output.txt}. Программа должна обрабатывать строки входного файла до его окончания. Каждая строка может иметь следующий формат:

\begin{itemize}
  \setlength{\itemsep}{0pt}
  \item \texttt{+ word 34} добавляет слово с номером 34. Ответ \texttt{OK} означает успех, \texttt{Exist} означает, что слово уже есть в словаре.
  \item \texttt{- word} удаляет слово. Ответ \texttt{OK} означает успех, \texttt{NoSuchWord} означает, что слово не найдено.
  \item \texttt{word} ищет слово. Для найденного слова выводится его номер, например \mbox{\texttt{OK: 34}}. Если слова нет, выводится \texttt{NoSuchWord}.
  \item \texttt{! Save dictionary.bin} сохраняет словарь в указанном файле в компактном бинарном представлении. \texttt{! Load dictionary.bin} загружает словарь и заменяет им текущий. При успехе выводится \texttt{OK}. При ошибке загрузки рабочий словарь должен остаться без изменений. Пустой или отсутствующий файл считается корректным и задаёт пустой словарь.
\end{itemize}

При системной ошибке, а также при ошибке сохранения или загрузки программа выводит \texttt{ERROR:} с описанием ошибки на английском языке. Это относится и к файлу неверного формата. Ограничения составляют 15 секунд и 512~МБ памяти.

Для реализации словаря необходимо провести исследование скорости выполнения и потребления оперативной памяти. В случае выявления ошибок или явных недочётов требуется их исправить.
\par
\endgroup
